\section{Power coefficients and their twists}
\label{sec:power-conventions-v158}
For a vector space $V$, put $M(V)=\Sym^2V$, with congruence action
$A\mapsto gAg^{\mathsf t}$, and define
\begin{equation}\label{eq:envelope-presentation-v158}
 B_h(V)=\Sym(M(V))/\langle (v^2)^{h+1}:v\in V\rangle .
\end{equation}
The degree of $M(V)$ is one. The generating space of the ideal is the
Cartan inclusion $\Sym^{2h+2}V\hookrightarrow\Sym^{h+1}(\Sym^2V)$,
normalized by $v^{2h+2}\mapsto(v^2)^{h+1}$. This definition requires
neither a volume form nor a complementary space. It agrees with
$B_{n,h}$ in Theorem~\ref{thm:quadratic-section-v155} when $W=V^*$.

\begin{proposition}[The coefficient map and its twists]
\label{prop:coefficient-conventions-v158}
Let $n=\dim V$ and $N=nh$. The quotient multiplication $\mu_p$ gives
an equivariant injection
\begin{equation}\label{eq:coefficient-duals-v158}
 \mu_p^*:B_h(V)_p^*\hookrightarrow\Sym^p(M(V)^*),
 \qquad (\mu_p^*\varphi)(A)=\varphi(A^p).
\end{equation}
Its image is the full coefficient space defining $J_p$, without a
choice of apolar coordinates. The socle and perfect pairings are
\begin{equation}\label{eq:perfect-twist-v158}
 B_h(V)_N=(\det V)^{2h},\qquad
 B_h(V)_p^*\simeq B_h(V)_{N-p}\otimes(\det V^*)^{2h}.
\end{equation}
For a line $L$, $B_h(V\otimes L)_p=B_h(V)_p\otimes L^{2p}$,
compatibly with all products. On $\Pj(M(V))$, the $p$th power is a
section of $B_h(V)_p\otimes\OO(p)$, and its base ideal is the image
of $B_h(V)_p^*\otimes\OO(-p)\to\OO$.
\end{proposition}
\begin{proof}
A homogeneous polynomial of degree $p$ is identified with its symmetric
$p$-linear polarization, normalized so that its value on $A$ is its
value on $(A,\ldots,A)$. Dualizing the surjection $\mu_p$ gives exactly
\eqref{eq:coefficient-duals-v158}. Multinomial coefficients appearing
in a monomial basis express this normalization and are nonzero; no
dual representation is silently replaced by the original one.
Theorem~\ref{thm:quadratic-section-v155} identifies the quotient with
determinant apolarity. The determinant of $A:V^*\to V$ takes values
in $(\det V)^2$. Choosing a trivialization gives a scalar polynomial
whose inverse-system line has character $(\det V^*)^{2h}$; its dual
socle therefore has character $(\det V)^{2h}$. The Gorenstein pairing
gives the second formula. Each generator in $M(V)$ has scalar weight
two, proving the line-twist assertion directly from
\eqref{eq:envelope-presentation-v158}. The tautological line
$\OO(-1)\to M(V)\otimes\OO$ has $p$th product
$\OO(-p)\to B_h(V)_p\otimes\OO$, which proves the last assertion.
\end{proof}


\subsection{Graph schemes and classical inputs}
The projective graph argument uses the following precise statement.
It applies to the ambient space and also to the two-dimensional total
base of the degeneration computed below.

\begin{lemma}[Simultaneous graph with invertible twists]
\label{lem:simultaneous-graph-v158}
Let $X$ be an integral noetherian complex scheme and let
$W_i\otimes L_i^{-1}\to\OO_X$ be nonzero coherent systems of
sections, with finite-dimensional $W_i$ and invertible $L_i$.
Write $K_i$ for their image ideals, and let
$U=X\setminus\bigcup_i V(K_i)$. The scheme-theoretic graph closure
of $U\to\prod_i\Pj(W_i^*)$ is
$\operatorname{Bl}_{\prod_iK_i}X$ with its specified graph embedding.
This statement identifies the scheme and the maps, not only their
normalizations. Replacing the product ideal by an invertible twist
or a positive power does not change its blow-up over $X$.
\end{lemma}
\begin{proof}
The open $U$ is dense because all $K_i$ are nonzero and $X$ is integral.
After the Segre embedding the simultaneous map is given by all products
of one section from each system. On a trivializing affine open, these
products generate $K=\prod_iK_i$. The homogeneous coordinate algebra
of the graph is the image of the polynomial algebra on these
products in $\OO_X[t]$, namely the Rees algebra
$\bigoplus_{m\geq0}K^m t^m$. It is a subalgebra of the function-field
polynomial ring. Thus its Proj is exactly the schematic closure of
the graph, with no extra component. The Segre equations hold on the
dense graph and hence on this integral closure scheme, so the graph
lands in the product, not merely in its ambient projective space.
Local trivializations glue. An invertible ideal changes the degree-$m$
piece by its $m$th tensor power and does not change relative Proj;
a positive power replaces the Rees algebra by a Veronese subalgebra.
\end{proof}

For the power system, $L_i=\OO(p)$ and $W_i=B_h(V)_p^*$;
for the exterior system of order $q$, $L_i=\OO(q)$ and its sections
are the $q$-minors. The determinant system is invertible on the
integral ambient space. It may be omitted from the blow-up product,
although its zero divisor must still be retained in contact formulas.
Likewise the first-coordinate identity map has unit base ideal.
These observations specify the projective twists in the proof of
Theorem~\ref{thm:power-graph-v157}.

The precise classical scheme-level complete-quadric construction is
\cite[Theorem~6.3]{Vainsencher1982}; the graph of exterior-power maps
and its successive symmetric-rank blow-ups are stated in
\cite[Remark~2.5 and Construction~2.6]{Massarenti2020}.
The relative rank version is recalled in
\cite[Definition~2.5, Remark~2.6 and Theorem~2.14]{CasarottiCornianiMassarenti2023}.
Consequently smoothness, the flag description, and the orbit
stratification in Corollary~\ref{cor:native-strata-v157} are classical
consequences of the graph identification, not new classifications.

There is also a precise invariant-ideal predecessor to the universal
identity. Henriques--Varbaro \cite[\S2.4, Theorems~2.6--2.8]{HenriquesVarbaro2016}
recall Abeasis's classification for the symmetric coordinate ring,
its products of determinantal ideals and their integral closures,
and the passage between invariant generators and products indexed
by diagrams. Our coefficient-space comparison uses this same
multiplicity-free setting. Theorem~\ref{thm:universal-power-v157}
identifies a particular multiplication coefficient map with the
balanced determinantal product as an ordinary ideal; it is not
an independent classification of invariant ideals. The exact
coefficient identification in Proposition~\ref{prop:coefficient-conventions-v158}
explains which map is being specialized on nonreduced bases.
We assert no exhaustive historical nonanticipation of the displayed
identity in all invariant-ideal terminology. Similarly the
multiplier calculation in Corollary~\ref{cor:power-multiplier-v157} is the
substitution of its exponents in the symmetric determinantal formula
\cite[Theorem~4.8]{HenriquesVarbaro2016}, not an independent
multiplier-ideal theorem.
